% kap02.tex — Lagrange-formalisme
% Hamilton-Jacobi-kompendiet

\chapter{Lagrange-formalisme}
\label{kap:lagrange}

\section{Målet: fra D'Alembert til én skalarfunktion}

I kapitel~\ref{kap:newton-variation} nåede vi til D'Alemberts princip,
\begin{equation}
  \sum_\alpha \left(\vb{F}_\alpha^{(a)} - m_\alpha\ddot{\vb{r}}_\alpha\right)\cdot\delta\vb{r}_\alpha = 0,
  \tag{\ref*{eq:dalembert}}
\end{equation}
som eliminerer tvangskræfterne, men stadig er udtrykt i de $3N$
kartesiske koordinater $\vb{r}_\alpha$ og deres virtuelle forskydninger
$\delta\vb{r}_\alpha$ — som ikke er uafhængige, når der er
tvangsbetingelser. Målet i dette kapitel er at omskrive princippet i de
$n$ uafhængige generaliserede koordinater $q_1,\dots,q_n$, hvor
$\delta q_1,\dots,\delta q_n$ \emph{er} uafhængige og arbitrære.
Resultatet bliver Euler-Lagrange-ligningerne — og undervejs vil vi se,
at hele bevægelsesligningen for et system kan pakkes ind i én eneste
skalarfunktion, Lagrange-funktionen $L$.

Vi udleder ligningerne på to uafhængige måder: først direkte fra
D'Alemberts princip (afsnit~\ref{sec:el-fra-dalembert}), og derefter fra
et rent variationsprincip, Hamiltons princip
(afsnit~\ref{sec:hamiltons-princip}). At de to veje fører til samme
ligninger, er ikke en tilfældighed — det er selve pointen, vi
identificerede i kapitel~\ref{kap:newton-variation}: det differentielle
og det globale synspunkt er ækvivalente udtryk for samme fysik.

\section{Euler-Lagrange-ligningerne fra D'Alembert}
\label{sec:el-fra-dalembert}

\subsection{Omskrivning af de virtuelle forskydninger}

Lad transformationsligningerne være $\vb{r}_\alpha =
\vb{r}_\alpha(q_1,\dots,q_n,t)$ som i \eqref{eq:transformationsligning}.
En virtuel forskydning (ved fast $t$, $\delta t=0$) hænger sammen med
de generaliserede virtuelle forskydninger $\delta q_i$ via kædereglen:
\begin{equation}
  \delta\vb{r}_\alpha = \sum_{i=1}^n \pdv{\vb{r}_\alpha}{q_i}\,\delta q_i.
  \label{eq:delta-r-kaederegel}
\end{equation}
Bemærk, at der \emph{ikke} optræder et $\partial\vb{r}_\alpha/\partial t$-led
her, netop fordi $\delta t=0$ per definition af en virtuel forskydning.

Vi indsætter \eqref{eq:delta-r-kaederegel} i D'Alemberts princip og får
to led, som vi behandler hver for sig.

\subsection{Det anvendte kraftled: generaliserede kræfter}

\begin{align}
  \sum_\alpha \vb{F}_\alpha^{(a)}\cdot\delta\vb{r}_\alpha
  &= \sum_\alpha \vb{F}_\alpha^{(a)}\cdot\sum_i \pdv{\vb{r}_\alpha}{q_i}\,\delta q_i \\
  &= \sum_i \left(\sum_\alpha \vb{F}_\alpha^{(a)}\cdot\pdv{\vb{r}_\alpha}{q_i}\right)\delta q_i
  \;\equiv\; \sum_i Q_i\,\delta q_i,
\end{align}
hvor vi har defineret den \emph{generaliserede
kraft}\index{generaliseret kraft}
\begin{equation}
  Q_i \equiv \sum_\alpha \vb{F}_\alpha^{(a)}\cdot\pdv{\vb{r}_\alpha}{q_i}.
  \label{eq:generaliseret-kraft}
\end{equation}
$Q_i$ behøver ikke have dimension af kraft — hvis $q_i$ er en vinkel,
er $Q_i$ et moment. Det er den generelle egenskab ved generaliserede
størrelser: $Q_i\,\delta q_i$ skal have dimension af arbejde, hvad
$Q_i$ og $\delta q_i$ hver for sig så end måler.

\subsection{Inertileddet}

Dette led er teknisk mere krævende:
\begin{equation}
  \sum_\alpha m_\alpha\ddot{\vb{r}}_\alpha\cdot\delta\vb{r}_\alpha
  = \sum_i \left(\sum_\alpha m_\alpha\ddot{\vb{r}}_\alpha\cdot\pdv{\vb{r}_\alpha}{q_i}\right)\delta q_i.
\end{equation}
Vi omskriver parentesen ved at bemærke identiteten
\begin{equation}
  \ddot{\vb{r}}_\alpha\cdot\pdv{\vb{r}_\alpha}{q_i}
  = \dv{t}\left(\dot{\vb{r}}_\alpha\cdot\pdv{\vb{r}_\alpha}{q_i}\right)
    - \dot{\vb{r}}_\alpha\cdot\dv{t}\left(\pdv{\vb{r}_\alpha}{q_i}\right),
  \label{eq:produktregel-identitet}
\end{equation}
som blot er produktreglen for differentiation, læst baglæns. Vi skal
nu bruge to hjælperesultater, der begge følger direkte af
$\vb{r}_\alpha=\vb{r}_\alpha(q_1,\dots,q_n,t)$:

\begin{itemize}
  \item \textbf{Hastighedsformlen.} Kædereglen giver
        \begin{equation}
          \dot{\vb{r}}_\alpha = \sum_j \pdv{\vb{r}_\alpha}{q_j}\dot q_j + \pdv{\vb{r}_\alpha}{t}.
          \label{eq:hastighed-kaederegel}
        \end{equation}
        Da $\vb{r}_\alpha$ ikke afhænger af $\dot q_j$, giver
        differentiation af \eqref{eq:hastighed-kaederegel} med hensyn
        til $\dot q_i$ direkte
        \begin{equation}
          \pdv{\dot{\vb{r}}_\alpha}{\dot q_i} = \pdv{\vb{r}_\alpha}{q_i}.
          \label{eq:hastighedsidentitet}
        \end{equation}
        Dette er den første ``kanoniske'' identitet, vi møder: den
        partielle afledte af \emph{hastigheden} med hensyn til en
        generaliseret \emph{hastighed} er den samme som den partielle
        afledte af \emph{positionen} med hensyn til den tilsvarende
        \emph{koordinat}.
  \item \textbf{Ombytning af afledte.} Da $\vb{r}_\alpha$ er en
        glat (to gange kontinuert differentiabel) funktion af
        $q_1,\dots,q_n,t$, kan vi ombytte den totale tidsafledte og den
        partielle afledte efter $q_i$:
        \begin{equation}
          \dv{t}\left(\pdv{\vb{r}_\alpha}{q_i}\right)
          = \sum_j \pdv{}{q_j}\left(\pdv{\vb{r}_\alpha}{q_i}\right)\dot q_j
            + \pdv{}{t}\left(\pdv{\vb{r}_\alpha}{q_i}\right)
          = \pdv{}{q_i}\left(\sum_j \pdv{\vb{r}_\alpha}{q_j}\dot q_j + \pdv{\vb{r}_\alpha}{t}\right)
          = \pdv{\dot{\vb{r}}_\alpha}{q_i},
          \label{eq:ombytning}
        \end{equation}
        hvor vi i sidste skridt igen brugte
        \eqref{eq:hastighed-kaederegel}.
\end{itemize}

Med \eqref{eq:hastighedsidentitet} og \eqref{eq:ombytning} bliver
\eqref{eq:produktregel-identitet}:
\begin{equation}
  \ddot{\vb{r}}_\alpha\cdot\pdv{\vb{r}_\alpha}{q_i}
  = \dv{t}\left(\dot{\vb{r}}_\alpha\cdot\pdv{\dot{\vb{r}}_\alpha}{\dot q_i}\right)
    - \dot{\vb{r}}_\alpha\cdot\pdv{\dot{\vb{r}}_\alpha}{q_i}.
\end{equation}
Vi ganger med $m_\alpha$, summerer over $\alpha$, og bemærker at
$m_\alpha\dot{\vb{r}}_\alpha\cdot\dot{\vb{r}}_\alpha = 2\cdot\frac12
m_\alpha \dot{\vb{r}}_\alpha^2$, så begge led kan skrives via den
kinetiske energi
\begin{equation}
  T \equiv \sum_\alpha \frac12 m_\alpha \dot{\vb{r}}_\alpha^2:
\end{equation}
\begin{equation}
  \sum_\alpha m_\alpha\ddot{\vb{r}}_\alpha\cdot\pdv{\vb{r}_\alpha}{q_i}
  = \dv{t}\pdv{T}{\dot q_i} - \pdv{T}{q_i}.
  \label{eq:inertiled-T}
\end{equation}

\subsection{Euler-Lagrange-ligningerne}

Vi samler de to led. D'Alemberts princip
$\sum_i\left(Q_i\,\delta q_i - \left[\dv*{}{t}\pdv*{T}{\dot q_i} - \pdv*{T}{q_i}\right]\delta q_i\right)=0$
gælder for \emph{alle} uafhængige, arbitrære $\delta q_i$ — og her er
den afgørende forskel fra kapitel~\ref{kap:newton-variation}: fordi
$q_i$ er valgt, så tvangsbetingelserne allerede er opfyldt, er
$\delta q_1,\dots,\delta q_n$ nu \emph{uafhængige}. Vi kan derfor sætte
hver koefficient til nul for sig:
\begin{equation}
  \boxed{\dv{t}\pdv{T}{\dot q_i} - \pdv{T}{q_i} = Q_i, \qquad i=1,\dots,n.}
  \label{eq:el-med-T-og-Q}
\end{equation}

Antag nu, at de anvendte kræfter er \emph{konservative}, dvs.\ at de
kan skrives som gradienten af en potentiel energi $V(\vb{r}_1,\dots,\vb{r}_N)$:
$\vb{F}_\alpha^{(a)} = -\grad_\alpha V$. Da er
\begin{equation}
  Q_i = \sum_\alpha \left(-\pdv{V}{\vb{r}_\alpha}\right)\cdot\pdv{\vb{r}_\alpha}{q_i}
      = -\pdv{V}{q_i},
\end{equation}
hvor $V$ nu forstås som funktion af $q_1,\dots,q_n,t$ via
\eqref{eq:transformationsligning}. Da $V$ ikke afhænger af
$\dot q_i$, kan vi flytte $-\partial V/\partial q_i$ ind under samme
afledte-struktur som $T$, ved at definere \emph{Lagrange-
funktionen}\index{Lagrange-funktion}
\begin{equation}
  \boxed{L(q,\dot q,t) \equiv T(q,\dot q,t) - V(q,t).}
  \label{eq:def-L}
\end{equation}
Da $\partial L/\partial \dot q_i = \partial T/\partial \dot q_i$ (fordi
$V$ ikke afhænger af $\dot q_i$), bliver \eqref{eq:el-med-T-og-Q} til
den berømte form:
\begin{equation}
  \boxed{\dv{t}\pdv{L}{\dot q_i} - \pdv{L}{q_i} = 0, \qquad i=1,\dots,n.}
  \label{eq:euler-lagrange}
\end{equation}
Dette er \emph{Euler-Lagrange-ligningerne}\index{Euler-Lagrange-ligningerne}. De er $n$ andenordens
differentialligninger — én for hver frihedsgrad — og de indeholder
\emph{al} den dynamiske information om systemet, uden at
tvangskræfterne nogen steder er optrådt eksplicit.

\begin{bemaerkning}
Det er værd at bemærke, hvor meget arbejde der er gemt i denne ene
ligning. Newtons anden lov skal opskrives separat for hver partikel, i
et koordinatsystem, hvor tvangskræfterne optræder eksplicit. Med
$L=T-V$ nedskrevet i de generaliserede koordinater, der passer til
problemets geometri, følger \emph{alle} bevægelsesligningerne
automatisk af \eqref{eq:euler-lagrange} — inklusive korrekte ``fiktive''
kræfter som centripetalacceleration, der ellers let glemmes eller
indføres ad hoc.
\end{bemaerkning}

\section{Eksempel: perlen på en roterende ring}

\begin{eksempel}
En perle med masse $m$ kan glide uden friktion langs en cirkulær ring
med radius $R$, som selv roterer med konstant vinkelhastighed $\omega$
om en lodret akse gennem ringens centrum (se
figur~\ref{fig:roterende-ring}). Perlens position på ringen beskrives
ved vinklen $\theta$ målt fra ringens laveste punkt; ringens egen
rotationsvinkel om den lodrette akse er $\phi=\omega t$, som er
\emph{foreskrevet}, ikke en dynamisk frihedsgrad. Systemet har derfor
kun \'en frihedsgrad, $q_1=\theta$.

I et fast (ikke-roterende) koordinatsystem med origo i ringens centrum
er perlens position
\begin{align}
  x &= R\sin\theta\cos(\omega t), \\
  y &= R\sin\theta\sin(\omega t), \\
  z &= -R\cos\theta.
\end{align}
Hastighedskomponenterne (brug produktreglen, og husk at $\theta=\theta(t)$):
\begin{align}
  \dot x &= R\dot\theta\cos\theta\cos(\omega t) - R\omega\sin\theta\sin(\omega t), \\
  \dot y &= R\dot\theta\cos\theta\sin(\omega t) + R\omega\sin\theta\cos(\omega t), \\
  \dot z &= R\dot\theta\sin\theta.
\end{align}
Den kinetiske energi $T=\frac12 m(\dot x^2+\dot y^2+\dot z^2)$
reduceres efter brug af $\cos^2+\sin^2=1$ (gennemfør selv
udregningen — det er en fin øvelse i at holde styr på krydsled) til
\begin{equation}
  T = \frac12 m R^2\dot\theta^2 + \frac12 m R^2\omega^2\sin^2\theta.
\end{equation}
Det andet led — som ikke involverer $\dot\theta$ — er netop den
energi, der i et roterende referencesystem ville vise sig som
centrifugalpotentialet. Her optræder det helt naturligt fra
transformationen, uden at vi nogensinde behøvede tale om
``centrifugalkraft''.

Den potentielle energi er $V = mgz = -mgR\cos\theta$, så
\begin{equation}
  L = T - V = \frac12 mR^2\dot\theta^2 + \frac12 mR^2\omega^2\sin^2\theta + mgR\cos\theta.
\end{equation}
Euler-Lagrange-ligningen \eqref{eq:euler-lagrange} for $\theta$ kræver
\begin{equation}
  \pdv{L}{\dot\theta} = mR^2\dot\theta, \qquad
  \dv{t}\pdv{L}{\dot\theta} = mR^2\ddot\theta, \qquad
  \pdv{L}{\theta} = mR^2\omega^2\sin\theta\cos\theta - mgR\sin\theta,
\end{equation}
og giver
\begin{equation}
  mR^2\ddot\theta - mR^2\omega^2\sin\theta\cos\theta + mgR\sin\theta = 0
  \quad\Longrightarrow\quad
  \ddot\theta = \omega^2\sin\theta\cos\theta - \frac{g}{R}\sin\theta.
\end{equation}
Bemærk ligevægtspunkterne: $\theta=0$ (perlen i bunden) er altid en
ligevægt, men den er kun \emph{stabil}, hvis $\omega^2 < g/R$. For
$\omega^2 > g/R$ bliver $\theta=0$ ustabil, og der opstår i stedet et
nyt, stabilt ligevægtspunkt ved $\cos\theta^\ast = g/(R\omega^2)$ — et
simpelt eksempel på en bifurkation, som ville have krævet omhyggelig
bogføring af den roterende ramme og centrifugalkraften, hvis vi havde
angrebet problemet med Newtons love direkte.
\end{eksempel}

\begin{figure}[htbp]
  \centering
  \begin{tikzpicture}[scale=1.0]
    % Lodret akse
    \draw[dashed] (0,-2.3) -- (0,2.3);
    % Ringen (set i perspektiv som ellipse)
    \draw[thick] (0,0) ellipse (1.8 and 0.55);
    \draw[thick,rotate around={90:(0,0)}] (0,0) ellipse (1.8 and 0.55);
    % Perlen
    \pgfmathsetmacro{\thetadeg}{55}
    \coordinate (B) at ({1.8*sin(\thetadeg)},{-1.8*cos(\thetadeg)});
    \filldraw[fill=blue!25,draw=blue!60!black,thick] (B) circle (0.13);
    % Vinkel theta fra bund
    \draw[dashed] (0,0) -- (0,-1.8);
    \draw[dashed] (0,0) -- (B);
    \draw (0,-0.7) arc ({-90+0}:{-90+\thetadeg}:0.7);
    \node at ({0.9*sin(\thetadeg/2+45)},{-0.9*cos(\thetadeg/2+45)}) {$\theta$};
    % Rotationspil for ringen om lodret akse
    \draw[-{Latex},thick,green!55!black] (1.4,1.1) arc (0:250:0.5);
    \node[green!55!black] at (1.4,1.75) {$\omega$};
  \end{tikzpicture}
  \caption{Perle (blå) på en cirkulær ring, der roterer med konstant
  vinkelhastighed $\omega$ om den lodrette akse. Perlens eneste
  frihedsgrad er vinklen $\theta$ målt fra ringens laveste punkt.}
  \label{fig:roterende-ring}
\end{figure}

\section{Hamiltons princip: et variationssyn}
\label{sec:hamiltons-princip}

Vi udleder nu Euler-Lagrange-ligningerne \eqref{eq:euler-lagrange} en
anden vej — ikke fra D'Alemberts differentielle princip, men fra et
globalt \emph{variationsprincip}, i tråd med udblikket i slutningen af
kapitel~\ref{kap:newton-variation}.

\subsection{Virkningen som funktional}

Betragt en tænkt bane $q_i(t)$ mellem to fastholdte tidspunkter $t_1$
og $t_2$, med fastholdte endepunkter
$q_i(t_1)=q_i^{(1)}$ og $q_i(t_2)=q_i^{(2)}$. For hver sådan bane
definerer vi \emph{virkningen}\index{virkning (Hamiltons principfunktion)}
\begin{equation}
  S[q] \equiv \int_{t_1}^{t_2} L(q(t),\dot q(t),t)\,dt.
  \label{eq:virkning}
\end{equation}
$S$ er en \emph{funktional}\index{funktional}: den tager en hel funktion $q(t)$ som
input og giver et enkelt tal som output — i modsætning til en
almindelig funktion, der tager et tal og giver et tal.
\emph{Hamiltons princip} siger:

\begin{center}
\emph{Den faktiske bane, systemet følger, er den bane, der gør $S$
\textbf{stationær} i forhold til alle nærliggende baner med samme
fastholdte endepunkter.}
\end{center}

``Stationær'' betyder her det direkte analogon til, at den afledte er
nul ved et ekstremum for en almindelig funktion — men nu i
funktionalers verden.

\subsection{Variation af virkningen}

Lad $q_i(t)$ være den (endnu ukendte) faktiske bane, og betragt en
nærliggende, tænkt bane
\begin{equation}
  \tilde q_i(t) = q_i(t) + \delta q_i(t),
\end{equation}
hvor $\delta q_i(t)$ er en arbitrær, lille afvigelsesfunktion, som
opfylder $\delta q_i(t_1)=\delta q_i(t_2)=0$ (endepunkterne er
fastholdt). Ændringen i virkningen til første orden i $\delta q_i$ er
\begin{equation}
  \delta S = \int_{t_1}^{t_2}\sum_i\left(\pdv{L}{q_i}\delta q_i + \pdv{L}{\dot q_i}\delta\dot q_i\right)dt,
\end{equation}
hvor $\delta\dot q_i = \dv*{}{t}\delta q_i$, fordi variation og
tidsafledning kommuterer (begge er blot grænseværdier af
differenskvotienter, og grænseovergangene kan ombyttes for glatte
funktioner). Vi integrerer det andet led partielt:
\begin{equation}
  \int_{t_1}^{t_2}\pdv{L}{\dot q_i}\dv{}{t}(\delta q_i)\,dt
  = \left[\pdv{L}{\dot q_i}\delta q_i\right]_{t_1}^{t_2}
    - \int_{t_1}^{t_2}\dv{}{t}\left(\pdv{L}{\dot q_i}\right)\delta q_i\,dt.
\end{equation}
Randtermen forsvinder, netop fordi $\delta q_i(t_1)=\delta q_i(t_2)=0$.
Vi står tilbage med
\begin{equation}
  \delta S = \int_{t_1}^{t_2}\sum_i\left[\pdv{L}{q_i} - \dv{}{t}\left(\pdv{L}{\dot q_i}\right)\right]\delta q_i\,dt.
  \label{eq:delta-S}
\end{equation}

\subsection{Variationsregningens grundlemma}

For at $\delta S=0$ skal gælde for \emph{enhver} tilladt afvigelse
$\delta q_i(t)$ (med de nævnte randbetingelser), må hver
koefficientfunktion i integranden være identisk nul. Dette følger af
\emph{variationsregningens fundamentallemma}: hvis en kontinuert
funktion $f(t)$ opfylder $\int_{t_1}^{t_2}f(t)\,\eta(t)\,dt=0$ for
\emph{alle} glatte $\eta(t)$ med $\eta(t_1)=\eta(t_2)=0$, så er
$f(t)\equiv 0$ på hele intervallet. (Beviset er enkelt: var $f(t_0)\neq0$
for et $t_0$ i det indre af intervallet, kunne man vælge et $\eta$, der
er positivt og koncentreret i en lille omegn om $t_0$ og nul udenfor —
og integralet ville da have samme tegn som $f(t_0)$, i modstrid med at
det skal være nul for alle $\eta$.)

Anvendt på \eqref{eq:delta-S}, koordinat for koordinat, giver dette
netop
\begin{equation}
  \dv{t}\pdv{L}{\dot q_i} - \pdv{L}{q_i} = 0, \qquad i=1,\dots,n,
\end{equation}
altså \emph{samme} Euler-Lagrange-ligninger \eqref{eq:euler-lagrange}
som i afsnit~\ref{sec:el-fra-dalembert} — nu udledt fra et rent
variationsprincip om hele banen, uden nogen reference til kræfter,
tvangskræfter eller virtuelle forskydninger ved et enkelt tidspunkt.

\begin{bemaerkning}
Det er denne ækvivalens, der gør Lagrange-formalismen så kraftfuld.
Vi har to fuldstændigt forskellige fysiske billeder — ``kræfterne
balancerer på hvert øjeblik'' (D'Alembert) og ``naturen vælger den bane,
der gør virkningen stationær'' (Hamilton) — der viser sig at give
\emph{identisk} matematisk indhold. Det er Hamiltons princip, ikke
D'Alemberts, der generaliserer til feltteori
(kapitel~\ref{kap:felter}) og til kvantemekanikkens
sti-integral-formulering, fordi det fra starten er formuleret som et
udsagn om en hel bane/historie snarere end om kræfter ved et punkt.
\end{bemaerkning}

\section{Tvangsbetingelser og Lagrange-multiplikatorer}

Indtil nu har vi forudsat, at tvangsbetingelserne allerede er
elimineret ved valg af generaliserede koordinater. Nogle gange er det
imidlertid nyttigt — eller nødvendigt, hvis man rent faktisk vil vide,
hvor stor en tvangskraft er (fx spændingen i en snor) — at beholde
flere koordinater end frihedsgrader og indføre tvangsbetingelserne
eksplicit.

Antag, at der er $m$ holonome tvangsbetingelser på formen
$f_k(q_1,\dots,q_N,t)=0$, $k=1,\dots,m$, blandt $N$ koordinater (hvor
$N>n=N-m$ er det faktiske antal frihedsgrader). Man kan da vise — ved
en variation af virkningen, hvor $\delta q_j$ nu er bundet af
$\sum_j (\partial f_k/\partial q_j)\delta q_j = 0$ for hvert $k$, og ved
at indføre Lagrange-multiplikatorer\index{Lagrange-multiplikatorer} $\lambda_k(t)$ til at håndtere
disse bindinger — at de udvidede Euler-Lagrange-ligninger bliver
\begin{equation}
  \dv{t}\pdv{L}{\dot q_j} - \pdv{L}{q_j} = \sum_{k=1}^m \lambda_k\,\pdv{f_k}{q_j}, \qquad j=1,\dots,N.
  \label{eq:el-multiplikator}
\end{equation}
De $m$ led på højresiden er netop de generaliserede tvangskræfter,
udtrykt via multiplikatorerne $\lambda_k$ — som selv bestemmes ved at
løse \eqref{eq:el-multiplikator} sammen med tvangsligningerne
$f_k=0$. Denne metode er praktisk, netop når man \emph{vil} kende
tvangskraftens størrelse (fx spændingen $\lambda$ i en snor), men vi
vil i resten af kompendiet næsten altid vælge koordinater, der
eliminerer tvangsbetingelserne fra start, som i
afsnit~\ref{sec:el-fra-dalembert}.

\section{Symmetrier og bevarelseslove: Noethers sætning}

Vi slutter kapitlet med et resultat, der viser Lagrange-formalismens
fulde styrke, og som bliver central senere i kompendiet: forbindelsen
mellem \emph{symmetrier} i $L$ og \emph{bevarelseslove}.

\subsection{Cykliske koordinater}

Antag, at $L$ ikke afhænger af en bestemt koordinat $q_k$ — dvs.
$\partial L/\partial q_k = 0$, selvom $L$ udmærket kan afhænge af
$\dot q_k$. En sådan koordinat kaldes \emph{cyklisk}\index{cyklisk koordinat}. Euler-Lagrange-
ligningen for $q_k$ reduceres da til
\begin{equation}
  \dv{t}\pdv{L}{\dot q_k} = 0
  \quad\Longrightarrow\quad
  p_k \equiv \pdv{L}{\dot q_k} = \text{konstant}.
  \label{eq:cyklisk-bevarelse}
\end{equation}
Størrelsen $p_k$ kaldes den \emph{generaliserede (kanoniske)
impuls}\index{kanonisk impuls}
konjugeret til $q_k$, og vi ser her det første eksempel på et
gennemgående mønster: \emph{en kontinuert symmetri i $L$ (uafhængighed
af en koordinat) giver en bevaret størrelse}. Vi vil møde $p_k$ igen
som en central variabel, når vi går videre til Hamilton-formalismen i
kapitel~\ref{kap:hamilton}.

\begin{eksempel}
For en fri partikel i et plan, $L=\frac12 m(\dot x^2+\dot y^2)$, er
både $x$ og $y$ cykliske. De bevarede impulser er
$p_x=m\dot x$ og $p_y=m\dot y$ — de sædvanlige kartesiske
bevægelsesmængder. For en central kraft i planpolar-koordinater,
$L=\frac12 m(\dot r^2+r^2\dot\phi^2)-V(r)$, er $\phi$ cyklisk (fordi
$V$ kun afhænger af $r$), og den bevarede størrelse er
$p_\phi=mr^2\dot\phi$ — det \emph{arealhastighedsmoment}, vi kender fra
Keplers anden lov.
\end{eksempel}

\subsection{Noethers sætning (translationssymmetri i tid)}

Cyklicitet er et specialtilfælde af en mere generel sætning, formuleret
af Emmy Noether\index{Noethers sætning}: \emph{til hver kontinuert symmetri af virkningen svarer
en bevaret størrelse}. Vi viser her det simpleste og for vores formål
vigtigste eksempel: symmetri under tidstranslation, som giver
energibevarelse.

Antag, at $L$ ikke afhænger eksplicit af $t$: $\partial L/\partial t=0$
(systemet er ``tidsinvariant'' — de fysiske love, det følger, ændrer sig
ikke fra et tidspunkt til et andet). Vi beregner den totale tidsafledte
af $L$ langs en faktisk løsning $q_i(t)$:
\begin{equation}
  \dv{L}{t} = \sum_i\left(\pdv{L}{q_i}\dot q_i + \pdv{L}{\dot q_i}\ddot q_i\right) + \pdv{L}{t}.
\end{equation}
Da $\partial L/\partial t=0$ per antagelse, og da vi langs en faktisk
løsning kan bruge Euler-Lagrange-ligningen
$\partial L/\partial q_i = \dv*{}{t}(\partial L/\partial \dot q_i)$ til
at omskrive det første led:
\begin{equation}
  \dv{L}{t} = \sum_i\left[\dv{}{t}\left(\pdv{L}{\dot q_i}\right)\dot q_i + \pdv{L}{\dot q_i}\ddot q_i\right]
  = \dv{}{t}\sum_i \dot q_i\pdv{L}{\dot q_i},
\end{equation}
hvor sidste skridt er produktreglen læst baglæns. Vi kan derfor skrive
\begin{equation}
  \dv{}{t}\left(\sum_i \dot q_i\pdv{L}{\dot q_i} - L\right) = 0,
\end{equation}
altså er størrelsen
\begin{equation}
  \boxed{H \equiv \sum_i \dot q_i\pdv{L}{\dot q_i} - L}
  \label{eq:def-H-forlobig}
\end{equation}
bevaret, når $L$ ikke afhænger eksplicit af tiden. Denne størrelse
kalder vi $H$ — og notationen er ikke tilfældig: dette er præcis den
kombination, der bliver til \emph{Hamilton-funktionen}, når vi i
kapitel~\ref{kap:hamilton} genindfører den systematisk via en
Legendre-transformation af $L$. For de fleste systemer, vi møder her —
hvor tvangsbetingelserne er tidsuafhængige, og $V$ ikke afhænger af
hastigheder — falder $H$ sammen med den samlede mekaniske energi
$T+V$, men vi udskyder det fulde bevis og den generelle diskussion
(inklusive tilfælde, hvor de \emph{ikke} falder sammen) til
kapitel~\ref{kap:hamilton}.

\begin{bemaerkning}
Vi har her set to eksempler på samme mønster: uafhængighed af $q_k$
(rumlig symmetri i en koordinat) giver bevaret impuls $p_k$;
uafhængighed af $t$ (tidssymmetri) giver bevaret ``energi'' $H$. Dette er
ikke en tilfældighed, men et specialtilfælde af Noethers generelle
sætning, som forbinder \emph{enhver} kontinuert symmetri af $L$ (under
en familie af transformationer parametriseret af én reel parameter)
med en tilsvarende bevaret størrelse. Den fulde, generelle formulering
af Noethers sætning — for en vilkårlig kontinuert transformation af
$q_i$ og $t$ — er ikke nødvendig for vores formål her, men mønstret
``symmetri $\leftrightarrow$ bevarelseslov'' vender vi tilbage til i
kapitel~\ref{kap:felter}, hvor det bliver et af de bærende principper
for klassisk feltteori.
\end{bemaerkning}

\section{Opsummering}

\begin{itemize}
  \item Ved at udtrykke D'Alemberts princip i generaliserede
        koordinater — hvor $\delta q_i$ er uafhængige — fås
        Euler-Lagrange-ligningerne
        $\dv*{}{t}(\partial L/\partial\dot q_i) - \partial L/\partial q_i = 0$,
        med $L=T-V$.
  \item Samme ligninger følger fra Hamiltons princip: den faktiske bane
        gør virkningen $S=\int L\,dt$ stationær blandt alle baner med
        fastholdte endepunkter. De to udledninger er ækvivalente, men
        det er variationssynet (Hamiltons princip), der generaliserer
        til feltteori og kvantemekanik.
  \item Tvangsbetingelser, man ikke ønsker at eliminere ved
        koordinatvalg, kan håndteres med Lagrange-multiplikatorer,
        som samtidig giver tvangskraftens størrelse.
  \item En cyklisk koordinat ($\partial L/\partial q_k=0$) giver en
        bevaret generaliseret impuls $p_k=\partial L/\partial\dot q_k$.
        Tidsuafhængighed af $L$ giver på samme måde en bevaret
        størrelse $H=\sum_i\dot q_i(\partial L/\partial\dot q_i)-L$ —
        et forvarsel om Hamilton-funktionen i
        kapitel~\ref{kap:hamilton}.
\end{itemize}

\begin{opgave}
Gennemfør selv den udregning af $T=\frac12 m(\dot x^2+\dot y^2+\dot z^2)$,
der blev sprunget over i eksemplet med den roterende ring, og
bekræft resultatet $T=\frac12 mR^2\dot\theta^2+\frac12 mR^2\omega^2\sin^2\theta$.
\end{opgave}

\begin{opgave}
En partikel med masse $m$ bevæger sig i planet under indflydelse af en
central kraft med potential $V(r)$. Opskriv Lagrange-funktionen i
planpolar-koordinater $(r,\phi)$, vis at $\phi$ er cyklisk, og udled
bevægelsesligningen for $r(t)$ ved brug af den bevarede impuls
$p_\phi=mr^2\dot\phi$ til at eliminere $\dot\phi$. Sammenlign med det
velkendte ``effektive potential'' $V_\text{eff}(r)=V(r)+p_\phi^2/(2mr^2)$.
\end{opgave}

\begin{opgave}
Vis, at for et system, hvor $\vb{r}_\alpha$ ikke afhænger eksplicit af
tiden (tidsuafhængige tvangsbetingelser, så $\partial\vb{r}_\alpha/\partial t=0$)
og hvor $V$ ikke afhænger af hastigheder, er $T$ en homogen
andengradsfunktion af $\dot q_i$, og brug Eulers sætning om homogene
funktioner ($\sum_i \dot q_i\,\partial T/\partial \dot q_i = 2T$) til at
vise, at $H$ defineret i \eqref{eq:def-H-forlobig} netop reducerer til
$H=T+V$ i dette tilfælde.
\end{opgave}
